% SHARD-ID: LB-00-OVERVIEW
% SHARD-TITLE: Overview -- the lattice infrared triangle
% SHARD-SOURCES: HANDOFF_MPS_SOFT_THEOREM.md; PRD.md; paper/skeleton.md;
%   HANDOFF.md; claims/CLAIMS.md (statuses quoted, not restated); CLAUDE.md
%   (laws L1--L10); labbook/WRITING-GUIDE.md
% SHARD-CLAIMS: none owned.  This shard states no theorem and therefore
%   carries no provenance block; every result it mentions is named in
%   English and cross-referenced to the section that owns it.

\section{Overview: the lattice infrared triangle}
\label{sec:overview}

\subsection{What the continuum triangle is}

In four-dimensional quantum field theory three families of statements that
look unrelated turn out to be one statement seen from three sides. The first
family is the \emph{soft theorems}: when one massless quantum in a scattering
process has its frequency taken to zero, the amplitude factorises into a
universal kinematic multiplier times the amplitude without that quantum. The
second is the \emph{asymptotic symmetries}: transformations that act
nontrivially on the data at null or spatial infinity, taken modulo those that
die off there, form a genuine symmetry group of the theory rather than a
gauge redundancy. The third is the \emph{memory effects}: after a burst of
radiation has passed, detectors do not return to their initial configuration
but retain a permanent, zero-frequency offset. The triangle is the assertion
that the soft theorem is the Ward identity of the asymptotic symmetry, and
that the memory is the Fourier-zero-mode content of the same Ward identity.
Each edge is a change of variables, not a new physical input.

Two features of that story matter for what follows. First, the triangle is,
in the continuum, a family of formal relations: the asymptotic symmetry
algebra is defined through boundary conditions at infinity whose function
spaces are delicate, and the soft factorisation is an all-orders statement
about an S-matrix that is not itself constructed. Second, there are two quite
different versions of it. The version usually advertised is the gauge one
(photons, gravitons), in which the soft factor has a $1/\omega$ pole because
a massless mediator is exchanged. The version relevant here is the
\emph{global} one: the soft-pion, or Goldstone, triangle, in which the soft
factor has no pole at all but an \emph{Adler zero} --- it vanishes linearly
as the momentum goes to zero, because the zero-momentum Goldstone mode is
exactly a global symmetry rotation of the vacuum and therefore acts trivially.

\subsection{Why the lattice version must be the Goldstone one}

This campaign asks whether the triangle is a theorem anywhere, and answers
the question on a one-dimensional quantum spin chain written in
matrix-product-state (MPS) language. There the three corners become finite,
checkable statements about tensors: the symmetry acts on-site, the vacua are
injective MPS, the excitations are ansatz vectors, and the Ward identity is
an exact algebraic rearrangement of a tensor network.

The gauge triangle has no home here. A one-dimensional lattice spin system
carries no massless gauge mediator --- the natural candidate, the Schwinger
model, confines and gives its photon a mass --- so there is no mechanism for
a $1/\omega$ pole, and any claim of one in this setting is simply wrong. The
correct antecedent is the global-symmetry triangle: the soft object is a
Goldstone-like magnon, the expected behaviour is an Adler zero, and the
leading soft datum is a \emph{slope} rather than a residue. Two further
structural facts fix the setting. Mermin--Wagner and Coleman forbid
continuous type-A symmetry breaking in a one-dimensional ground state, so the
triangle must live either on a type-B Goldstone (the isotropic Heisenberg
ferromagnet, exact product vacuum, magnon dispersion $\omega\sim k^2$, exact
Bethe scattering data available as an independent oracle) or on a discretely
broken symmetry (the easy-axis XXZ ferromagnet, two product vacua, domain
walls, gapped magnons, so that ``soft'' means a low-frequency expansion about
the gap). And there are no boosts on a lattice, so the universal factor
cannot take a Lorentz-covariant form; its correct dependence on the legs'
group velocities is part of the problem rather than an obstacle to it.

The two model classes are used throughout. The ferromagnetic chain supplies
exact magnons and an exact two-body scattering matrix; the easy-axis chain
supplies exact kinks and a wall-position observable whose displacement is the
memory. Neither model alone exhibits every corner, and that tension is
structural: it is respected in what follows, not papered over.

\subsection{The four headline findings}

\paragraph{(i) Memory quantization needs no scattering theory.}
The naive expectation, inherited from the continuum, is that the memory is
quantized because it is the zero-frequency limit of a soft factor, so that
one must first have a soft theorem and an S-matrix before one can say
anything about the permanent displacement. That expectation is wrong here,
and its failure is the campaign's most robust positive result. Fix a finite
window $W=[a,b]$ of the chain, a cut $c_0\in W$, and the windowed
wall-position observable
\begin{equation}
  \wallX_W \;=\; a-1+\frac{1}{2s}\sum_{x=a}^{b}\bigl(S^z_x+s\bigr),
  \qquad
  \Qhat_{W,c_0} \;=\; 2s\bigl(\wallX_W-c_0\bigr),
  \label{eq:ov-window-charge}
\end{equation}
where $\pm s$ are the charge densities of the two vacua. If the on-site
charge merely generates a circle --- that is, if $e^{2\pi i S^z_x}$ is a
scalar, which holds for every spin-$S$ chain --- then the spectrum of
$\Qhat_{W,c_0}$ lies in a single coset of $\Z$, and the increment recorded by
measuring this same observable before and after an event is an
\emph{integer}. No wave operators, no channel inventory, no completeness, no
gap, no Lieb--Robinson velocity: the offsets at the two measurement times are
the same number and cancel. This is the finite-window memory-index theorem
(\cref{sec:memory-index}), and it is unconditional. Its infinite-volume
companion --- that every ordered limit law of the measurement protocol is a
probability distribution supported on $\Z$, with
$\delta x=-(2s)^{-1}\sum_\nu \nu\,p_\nu$ --- is proved as a conditional
implication whose hypotheses are relaxation hypotheses on the charge history,
never scattering hypotheses; and one of those hypotheses turns out to be a
free theorem, provable from separability of the algebra and strong continuity
of the dynamics alone (\cref{sec:local-relaxation}). Symmetry, not dynamics,
does the quantizing.

\paragraph{(ii) Scattering only computes the value.}
What the quantized ledger does not supply is the number multiplying it. That
number is on-shell data, and it is obtained by honest two-body scattering.
For the fully polarised bilinear ferromagnet at every site spin
$S\in\{\tfrac12,1,\tfrac32,\dots\}$ the exact two-magnon ratio in the regular
channel is computed directly from the separated, double-occupancy and
adjacent equations, with no integrability or completeness assumed, and the
physical phase has the exact soft slope
\begin{equation}
  \left.\partial_{k_s}\delta_{\mathrm{phys}}\right|_{k_s=0}
  \;=\;\frac{\sgn(v_h-v_s)}{S}.
  \label{eq:ov-slope}
\end{equation}
This is the exact spin-$S$ two-magnon slope theorem (\cref{sec:two-magnon}),
\statusProved. The same value reaches the protocol datum only through an
identification step that is itself proved on a displayed
separated-preparation subclass (\cref{sec:soft-index}): the value is
transported into the infrared statement by matching, never installed by
stipulation. The division of labour is the campaign's spine sentence:
\emph{symmetry quantizes the infrared data; the excitation ansatz supplies
the kinematics; dynamics only picks the values.}

\paragraph{(iii) The naive universality conjecture is refuted.}
The brief's Conjecture~M --- that the wall displacement equals the
zero-frequency limit of the soft factor, a lattice Braginsky--Thorne relation
--- is \statusRefuted{} as literally stated (\cref{sec:memory-quantization}).
What survives is weaker and sharper: an exact identity expressing the
displacement as the finite-time direct-current weight of the \emph{physical}
boundary current, and a conditional charge-bookkeeping law. Independently,
the stronger claim that the soft factorisation is universal over unrestricted
local or quasi-local sources is also \statusRefuted{}
(\cref{sec:universality}): an explicit four-site source leaves the one-hard
amplitude untouched while shifting the linear soft coefficient. Universality
survives only relative to a named class of sources with no contact term, and
even there the class constant is left open by the factorisation argument
itself. These refutations are presented here as results, with the surviving
weaker statements attached, because a sharp refutation plus its survivor is
worth more than a pile of analogies.

\paragraph{(iv) An SPT rigidity dichotomy.}
When the symmetry is realised projectively on the virtual space --- the
one-dimensional symmetry-protected-topological (SPT) situation --- the
campaign separates two kinds of infrared coefficient
(\cref{sec:spt}). Registered \emph{edge} residues are rigid: in a fixed
Schmidt/edge register the compensated group residue is the projective
representation itself, the Hermitian residue has spectrum congruent to a
fixed offset modulo $\Z$, and its irreducible dimension is bounded below by
the minimal dimension of a projective irreducible representation of the
class. Normalised \emph{bulk} coefficients are not rigid: along a path with a
common gap and continuous external data they vary continuously, and become
topological only after a separate local-constancy argument that is not
supplied. Quantized edge residue, drifting bulk coefficient --- that
dichotomy, rather than a blanket ``the class appears in all soft data'', is
what the analysis supports.

\subsection{How to read this document}

Every result in this labbook is stated under a descriptive English name and
carries a coloured status tag. The five tags, and exactly what each asserts,
are:

\begin{description}
  \item[\statusProved] a structured proof exists and survived the capped
    hostile review round with no open fatal or major objection on the
    promoted statement.
  \item[\statusProvedConditional] the same, but what is proved is an
    implication. The antecedent is displayed inside the statement; it is
    never moved into a footnote and never assumed away.
  \item[\statusSketch] an argument exists and its gaps are named. The
    statement says precisely what is established and what is missing.
  \item[\statusConjecture] a precise statement with no proof. Whatever
    evidence exists --- a numerical falsifier that was survived, an oracle
    agreement --- is described as evidence and not as proof.
  \item[\statusRefuted] false as stated. The mechanism that kills it is
    given, together with the weaker statement that survives.
\end{description}

Immediately after each definition, theorem, conjecture or refutation a small
grey \emph{provenance block} records four things: the campaign identifier
used in the repository's claims file, the file and proof step where the
statement is established, the checker or numerical experiment that tests it
(or a dash, if none does), and the review record. That block is the only
place in this document where campaign identifiers carry weight. The prose
never argues from an identifier; it argues from named results and
cross-references.

Because the campaign has accumulated a large private vocabulary ---
claim identifiers, numbered definitions, named hypotheses in parentheses ---
\cref{sec:dictionary} is a complete dictionary. Every identifier that appears
anywhere in the campaign's claim ledger is listed there with a plain-English
gloss and a pointer to the section that owns it. A reader who meets an
unfamiliar tag in a provenance block should look it up there and nowhere
else; in particular, this document is meant to be readable front to back
without ever opening the repository.

The remaining sections follow the logical order rather than the historical
one. \Cref{sec:mps-setting,sec:symmetry-noether} fix the setting and the
lattice Noether pair; \cref{sec:corner-a} is the asymptotic-symmetry corner;
\cref{sec:kink-model,sec:memory-quantization,sec:memory-index,sec:local-relaxation}
build the memory corner, from the kink model through the quantization theorem
to the relaxation hypotheses it rests on;
\cref{sec:two-magnon,sec:wave-operators,sec:soft-index,sec:universality} build
the soft corner, from exact two-body data through wave operators to the soft
index and the universality class; \cref{sec:spt} is the rigidity dichotomy;
\cref{sec:two-plus-one} records what generalises to two spatial dimensions;
\cref{sec:negative-results} collects the refutations and obstructions in one
place; \cref{sec:numerics} the numerical instruments; and
\cref{sec:triangle-status} states, corner by corner and edge by edge, exactly
where the triangle stands. \Cref{sec:continuum-reduction} then asks whether all
of this reduces to what the continuum already accepts, and
\cref{sec:operational-status} asks whether any of it can be measured.

That last question is asked under a rule worth stating here, at the front,
because it governs how every object in this document should be read: \textbf{a
representation is never operationally meaningful.} Tensors, canonical forms,
virtual bond spaces, intertwiners and ansatz gauge data are bookkeeping; only
what some describable quantum experiment can measure, or some describable
quantum operation can implement, has operational content --- and each such
object must say which of the two it is, with the experiment or the invariant
named. \Cref{sec:operational-status} carries out that audit on every definition
and every load-bearing object in the campaign, gives six explicit protocols for
the quantities that pass, proves that one much-used ansatz coordinate is exactly
unmeasurable, and records what the operator-algebraic literature on boundary
observables already owns. A reader who wants to know whether a given symbol in
this document is physics or scaffolding should look there first.

\subsection{The honesty contract}

One rule governs this document above all others: \textbf{every status written
here matches the campaign's claim ledger exactly}. A conditional theorem is
labelled conditional and displays its condition. A statement whose proof
holds only on a restricted register, a restricted packet class, a projected
component, or a single model says so in the statement itself, not in a
remark. Nothing is rounded up: a sketch is never described as ``essentially
proved'', a survived falsifier is never described as a verification, and a
theorem about an ansatz-conditioned construction is never described as a
theorem about the true model. Where a claim was once stated more strongly and
has since been narrowed --- and there are several such cases, including one
projection identity that is exact for a single magnon and false for two ---
the erratum is written out rather than quietly absorbed.

The reason for the rule is not decorum. The campaign's one non-negotiable
constraint is that it must never carry a label stronger than its evidence: a
claim resting at sketch or conjecture indefinitely costs nothing, whereas a
single mislabelled statement would make the whole ledger worthless. This
labbook inherits that constraint, and any commit that changes a status in the
ledger is required to change the corresponding statement here in the same
breath.
